\chapter{Clouds, tests, and the 2-to-2 game}\label{ch:inner}

The smooth game is a projection game, not a unique game. On an edge $(U,O)$, each replaced equation makes the projection 2-to-1 in its entry, as in Figure~\ref{fig:evv}, so the projection is many-to-one unless no equation was replaced. How can its labels be checked by constraints between single elements of a small alphabet $K$? OpenAI's construction replaces each first endpoint of the smooth game by a set of new vertices, called its cloud, whose labelings encode labels of the endpoint, and it puts constraints inside each cloud that encodings of labels satisfy. For Part~I, the answer is the 2-to-2 game $G_{2:2}$ of Definition~\ref{def:G22}, whose value is at least $1-k\xi-2^{-\ell}-2^{1-m}$ when $\OPT(\Theta)\ge1-\xi$.

\paragraph{Clouds and encodings.}\phantomsection\label{nt:clouds} Fix $\ell\ge1$. A \emph{table} on a vertex $Q$ of the smooth game is a linear map $P:L_Q\to\F_2^\ell$, and the \emph{encoding} of a label $z\in L_Q$ is the function $P\mapsto Pz$ on tables. On a first endpoint $U$, the tables are the matrices $M\in\Mat_{\ell,m}=\Hom(L_U,\F_2^\ell)$, and the encoding of $z$ is $f_z(M)=Mz$. Each table on a first endpoint $U$ gives a vertex of a new game, with labels in $K=\F_2^\ell$. Tables on different first endpoints can give the same vertex, as the next paragraph explains. The tables on $U$ form the \emph{cloud} of $U$, and the encoding $f_z$ of a label $z$ of $U$ is a labeling of this cloud (Figure~\ref{fig:layers}).

\paragraph{Shared vertices.}\phantomsection\label{nt:shared} Clouds are joined through the second endpoints. A table $P'$ on $O$ gives the table $P'\bar\pi=P'\circ\bar\pi$ on $U$, its \emph{pullback} along $\bar\pi$. For every $z\in L_U$, $(P'\bar\pi)z=P'(\bar\pi z)$: the encoding of $z$ at the pullback equals the encoding of the restriction $\bar\pi z$ at $P'$. In OpenAI's reduction, a table $P'$ on $O$ and its pullbacks to all first endpoints $U$ with an edge $(U,O)$ are a single vertex\footnote{OpenAI's vertices are $K$-orbits of \emph{keys}, where the key of a table is a reduced form of it that keeps, in each entry, only the coordinates on which the table depends \cite[Section~4.3]{OAI}. A table on $O$ and its pullbacks have the same key, and other tables can share a key too.}, shared by the clouds of these $U$. Suppose the clouds of $U_1$ and $U_2$ are labeled by the encodings of $z_{U_1}$ and $z_{U_2}$, and let $\bar\pi_1$ and $\bar\pi_2$ be the restrictions from $U_1$ and $U_2$ to $O$. The shared vertex $P'$ then receives the label $P'(\bar\pi_1z_{U_1})$ from one cloud and $P'(\bar\pi_2z_{U_2})$ from the other. These agree for every table $P'$ on $O$ exactly when $\bar\pi_1z_{U_1}=\bar\pi_2z_{U_2}$, that is, when $z_{U_1}$ and $z_{U_2}$ restrict to the same label of $O$. So encodings $f_{z_U}$ of labels $z_U$, one on each cloud, take the same value at all pullbacks of each table on a second endpoint exactly when the labels $z_U$ come from a labeling of the smooth game that satisfies every edge. Hence the names: the smooth game is the \emph{outer game}, which determines the vertices that clouds share, and the constraints inside the clouds form the \emph{inner test}.

\begin{figure}[ht]
\centering
\begin{tikzpicture}[font=\small,
  dot/.style={circle,fill,inner sep=1.3pt},
  qv/.style={draw,circle,inner sep=1pt,minimum size=7mm}]
\node[qv] (u1) at (0.2,2.7) {$U_1$};
\node[qv] (o) at (1.9,2.7) {$O$};
\node[qv] (u2) at (3.6,2.7) {$U_2$};
\draw[-{Stealth[length=2mm]}] (u1) -- (o);
\draw[-{Stealth[length=2mm]}] (u2) -- (o);
\node[right,align=left] at (5.9,2.7) {outer game:\\edges of the smooth game};
\draw (0.2,0) ellipse (1.85 and 0.95);
\draw (3.6,0) ellipse (1.85 and 0.95);
\node at (-0.6,1.2) {cloud of $U_1$};
\node at (4.4,1.2) {cloud of $U_2$};
\node[dot] (a1) at (-0.9,0.35) {};
\node[dot] (a2) at (-0.4,-0.45) {};
\node[dot] (a3) at (0.5,0.5) {};
\node[dot] (s1) at (1.9,0.3) {};
\node[dot] (s2) at (1.9,-0.3) {};
\node[dot] (b1) at (3.3,0.5) {};
\node[dot] (b2) at (3.8,-0.45) {};
\node[dot] (b3) at (4.6,0.3) {};
\draw (a1) -- (a2); \draw (a2) -- (s2); \draw (a3) -- (s1);
\draw (s1) -- (b1); \draw (s2) -- (b2); \draw (b2) -- (b3);
\node[font=\footnotesize,align=center] (sl) at (1.9,-1.55) {shared vertices: a table $P'$ on $O$\\and its pullbacks $P'\bar\pi$};
\draw[-{Stealth[length=1.6mm]}] (sl) -- (s2);
\node[right,align=left] at (5.9,0) {new game: tables are its vertices,\\and a table on $O$ and its\\pullbacks count as one vertex.\\Inner-test edges join\\tables of one cloud.};
\end{tikzpicture}
\caption{The two layers in OpenAI's reduction (schematic). The edges of the outer game (top) determine which vertices the clouds of $U_1$ and $U_2$ share (bottom).}\label{fig:layers}
\end{figure}

Each inner test in these notes compares the labels at two tables $M$ and $M+al^\top$ of one cloud, where $a$ and $l$ are random vectors, so the two tables differ by a matrix of rank at most one.

\begin{definition}[{2-to-2 test: Barak, Kothari, and Steurer 2018 \cite[Section~1.3, Test~2]{BKS18}; equality test: \cite[Section~1.4]{OAI}}]\label{def:shortcode}
The encoding $f_z(M)=Mz$ of $z\in\F_2^m$ is its \emph{degree-2 shortcode}. For $f:\Mat_{\ell,m}\to\F_2^\ell$ and independent uniform $M\in\Mat_{\ell,m}$, $a\in\F_2^\ell$, $l\in\F_2^m$, the \emph{2-to-2 test} accepts if $f(M+al^\top)-f(M)\in\{0,a\}$, and the \emph{equality test} accepts if $f(M+al^\top)=f(M)$.
\end{definition}

Encodings pass the two tests with different probabilities.

\begin{lemma}\label{lem:completeness}
For $z\ne0$, $f_z$ passes the 2-to-2 test with probability $1$ and the equality test with probability $(1+2^{-\ell})/2$.
\end{lemma}

\begin{proof}
By linearity, $f_z(M+al^\top)-f_z(M)=a(l^\top z)\in\{0,a\}$, so the 2-to-2 test always accepts. The equality test accepts exactly when $l^\top z=0$ or $a=0$. Since $z\ne0$, the bit $l^\top z$ is uniform and independent of $a$, so this event has probability $\tfrac12+\tfrac12\cdot2^{-\ell}$.
\end{proof}

So encodings pass the 2-to-2 test with probability $1$, but its constraint is not a permutation. The equality test imposes a permutation, but encodings pass it with probability only $(1+2^{-\ell})/2$. Part~I uses the 2-to-2 test and obtains 2-to-2 games. Part~II replaces it by a test whose constraints are permutations and which encodings pass with probability close to $1$. For $a\ne0$, the 2-to-2 test imposes a 2-to-2 constraint (Chapter~\ref{ch:ug}) on $\bigl(f(M),f(M+al^\top)\bigr)$, whose blocks $B_j=B'_j$ are the $|K|/2$ cosets $\{y,y+a\}$ of $\{0,a\}$ (Figure~\ref{fig:twotwo}). Splitting it by Lemma~\ref{lem:split} gives the identity $\sigma_0$, which is the constraint of the equality test, and the translation $\sigma_1(y)=y+a$. The encoding $f_z$ satisfies $\sigma_0$ exactly when $a(l^\top z)=0$, and $\sigma_1$ exactly when $a(l^\top z)=a$. Given $a\ne0$, each event has probability exactly $1/2$, so splitting halves the completeness, as in Corollary~\ref{cor:half}.

\begin{figure}[ht]
\centering
\begin{tikzpicture}[font=\small,
  lb/.style={inner sep=2pt,minimum width=1.3cm}]
\node[lb] (l1) at (0,3) {$y$};
\node[lb] (l2) at (0,2.2) {$y+a$};
\node[lb] (l3) at (0,1) {$\tilde y$};
\node[lb] (l4) at (0,0.2) {$\tilde y+a$};
\node[lb] (r1) at (4,3) {$y$};
\node[lb] (r2) at (4,2.2) {$y+a$};
\node[lb] (r3) at (4,1) {$\tilde y$};
\node[lb] (r4) at (4,0.2) {$\tilde y+a$};
\foreach \n in {l,r} {
  \node[draw,rounded corners=3pt,fit=(\n1)(\n2),inner sep=1pt] {};
  \node[draw,rounded corners=3pt,fit=(\n3)(\n4),inner sep=1pt] {};
}
\node[above=0.45cm] at (l1) {label at $M$};
\node[above=0.45cm] at (r1) {label at $M+al^\top$};
\foreach \i in {1,2,3,4} { \draw[thick] (l\i.east) -- (r\i.west); }
\draw[dashed] (l1.east) -- (r2.west);
\draw[dashed] (l2.east) -- (r1.west);
\draw[dashed] (l3.east) -- (r4.west);
\draw[dashed] (l4.east) -- (r3.west);
\node at (0,-0.55) {$\vdots$};
\node at (4,-0.55) {$\vdots$};
\draw[thick] (6.0,2.2) -- (6.8,2.2) node[right] {$\sigma_0=$ identity: the equality test};
\draw[dashed] (6.0,1.2) -- (6.8,1.2) node[right] {$\sigma_1=$ translation by $a$};
\end{tikzpicture}
\caption{For $a\ne0$, the constraint of the 2-to-2 test is the union of the graphs of two permutations (Lemma~\ref{lem:split}). Two of the $|K|/2$ blocks are shown. Indexing each block $\{y,y+a\}$ by $y_0=y$ and $y_1=y+a$ gives $\sigma_c(y)=y+ca$. The pair $\bigl(f_z(M),f_z(M+al^\top)\bigr)$ lies on the graph of $\sigma_{l^\top z}$: on a solid edge when $l^\top z=0$ and on a dashed edge when $l^\top z=1$.}\label{fig:twotwo}
\end{figure}

\paragraph{Folding.}\phantomsection\label{nt:folding} A constant labeling passes both tests on every edge, since $f(M+al^\top)-f(M)=0$. So a game built from either test has value $1$ and cannot be sound. Folding excludes constant labelings. Chapter~\ref{ch:twotwo} uses two further consequences of folding: the vectors that its decoding returns have $z_0=1$ (Lemma~\ref{lem:labelset} and the $O$-prover's strategy), and a restored function, defined below, passes the 2-to-2 test with at most twice the probability with which it passes the equality test (Lemma~\ref{lem:foldcompare}). Write $P\in\Hom(L_Q,K)$ for a table on a vertex $Q$. For $h\in K$, let $P^{+h}=P+he_0^\top$ be $P$ with $h$ added to its column at coordinate $0$, where $e_0$ is the basis vector of the homogeneous coordinate. So $P^{+h}z=Pz+hz_0$. Folding makes each orbit $\{P^{+h}:h\in K\}$ one vertex $v$ with one label $\lambda(v)$, assigned to a representative $R_v$. It then defines the \emph{restored} function $F(R_v^{+h})=\lambda(v)+h$ (Figure~\ref{fig:folding}), which satisfies $F(P^{+h})=F(P)+h$ and so is not constant. For example, the equality test's constraint $F(R_v^{+h})=F(R_{v'}^{+h'})$ becomes the translation constraint $\lambda(v')=\lambda(v)+(h-h')$. Encodings satisfy the same identity: if $z_0=1$, then $f_z(M^{+h})=f_z(M)+h$. For a vertex $Q$ of the smooth game, write $F_Q$ for the restored function on the tables of $Q$. Folding and the shared vertices give the identities
\[
F_Q(P^{+h})=F_Q(P)+h\quad\text{(folding)},\qquad F_U(P'\bar\pi)=F_O(P')\quad\text{(sharing)},
\]
for every vertex $Q$, every edge $(U,O)$, and every table $P'$ on $O$. The soundness proofs of Chapters~\ref{ch:twotwo} and~\ref{ch:openai} use a labeling only through these identities and through the probability that its restored functions pass a test \cite[Proposition~5.3]{OAI}.

\begin{figure}[ht]
\centering
\begin{tikzpicture}[font=\small,t/.style={draw,minimum height=6mm,inner sep=3pt}]
\node[t] (p0) at (0,1.9) {$R_v$};
\node[t] (p1) at (2.7,1.9) {$R_v^{+h}$};
\node (pd) at (4.7,1.9) {$\cdots$};
\node[font=\footnotesize,above=1mm of p0] {$\lambda(v)$};
\node[font=\footnotesize,above=1mm of p1] {$\lambda(v)+h$};
\node[t] (q0) at (0,0) {$R_{v'}$};
\node[t] (q1) at (2.7,0) {$R_{v'}^{+h'}$};
\node (qd) at (4.7,0) {$\cdots$};
\node[font=\footnotesize,below=1mm of q0] {$\lambda(v')$};
\node[font=\footnotesize,below=1mm of q1] {$\lambda(v')+h'$};
\node[draw,dashed,rounded corners,fit=(p0)(pd),inner sep=5pt] (ov) {};
\node[draw,dashed,rounded corners,fit=(q0)(qd),inner sep=5pt] (ow) {};
\node[left=1mm of ov] {vertex $v$};
\node[left=1mm of ow] {vertex $v'$};
\draw[thick] (p1) -- node[right,align=left,font=\footnotesize] {test $F(R_v^{+h})=F(R_{v'}^{+h'})$, i.e.\\$\lambda(v')=\lambda(v)+(h-h')$} (q1);
\end{tikzpicture}
\caption{Folding: each orbit of tables is one vertex with one label. The equality test between two tables becomes a translation constraint between their labels.}\label{fig:folding}
\end{figure}

\paragraph{The 2-to-2 game.} Part~I's reduction maps a parity system $\Theta$ with three distinct variables in each equation (Chapter~\ref{ch:outer}) to a 2-to-2 game $G_{2:2}$. Its vertices are the $K$-orbits of keys (the footnote in ``Shared vertices''), with labels in $K$, and a labeling $\lambda$ gives restored functions $F_Q$ that satisfy the folding and sharing identities.

\begin{definition}[The 2-to-2 game]\label{def:G22}
The random edge of $G_{2:2}$ is drawn as follows. Sample $U=(\theta_1,\dots,\theta_k)$ with independent $\theta_i\sim p$, and independent uniform $M\in\Mat_{\ell,m}$, $a\in K$, and $l\in\F_2^m$. The edge joins the vertices of the tables $M$ and $M+al^\top$ on $U$, and its constraint is
\[
F_U(M+al^\top)-F_U(M)\in\{0,a\}.
\]
\end{definition}

So $\val_\lambda(G_{2:2})$ is the 2-to-2 test's acceptance probability for $F_U$, averaged over $U$ (Definition~\ref{def:shortcode}). With $M=R_v^{+h}$ and $M+al^\top=R_{v'}^{+h'}$, the constraint reads $\lambda(v')\in\lambda(v)+h+h'+\{0,a\}$. For $a\ne0$ this is a 2-to-2 constraint: its blocks for $v$ are the $|K|/2$ cosets of $\{0,a\}$, and its blocks for $v'$ are their translates by $h+h'$. On a fixed $U$, distinct tables have distinct keys \cite[Section~4.3]{OAI}. So the two tables lie in one vertex exactly when $al^\top=he_0^\top$ for some $h\in K$, that is, when $a=0$, $l=0$, or $l=e_0$. Every labeling satisfies these loops, by folding when $l=e_0$. Deleting these loops, which have total probability $\rho_0\le2^{-\ell}+2^{1-m}$, and renormalizing changes each value $\val_\lambda$ to a number in $[\val_\lambda-\rho_0,\,\val_\lambda]$. Since $k$ and $\ell$ are fixed, $G_{2:2}$ has size polynomial in that of $\Theta$.

\paragraph{Completeness.} Suppose an assignment $\mathbf x$ satisfies $\Theta$ with probability at least $1-\xi$. For a vertex $U$ whose equations $\mathbf x$ all satisfies, let $z_U\in L_U$ be the label of $U$ given by $\mathbf x$. OpenAI labels each vertex of $G_{2:2}$, a $K$-orbit of keys, at its representative: the label is the representative key evaluated at bits read from $\mathbf x$, namely the bit of $\mathbf x$ at each variable that the key keeps, and the first two bits of $\mathbf x$ at each equation that it keeps, with the third bit set by the equation. The restored functions of this labeling satisfy $F_U(M)=Mz_U$ for every such $U$ \cite[Section~4.3 and proof of Lemma~4.3]{OAI}, so $F_U$ passes the 2-to-2 test with probability $1$ (Lemma~\ref{lem:completeness}). The equations $\theta_1,\dots,\theta_k$ are independent with law $p$, so by a union bound $\mathbf x$ satisfies all of them with probability at least $1-k\xi$. Hence $\val(G_{2:2})\ge1-k\xi-\rho_0$ after deleting loops.

\chapterquestions
\begin{question}\label{q:homog}
Folding adds $h$ to the column of a table at the homogeneous coordinate. Why that column? What happens to the encodings if $h$ is added to the column at a coordinate $i\ne0$ instead?
\end{question}

\begin{question}\label{q:disagree}
Let $(U_1,O)$ and $(U_2,O)$ be edges of the smooth game, with restrictions $\bar\pi_1$ and $\bar\pi_2$, and let $z_1\in L_{U_1}$ and $z_2\in L_{U_2}$ be labels with $\bar\pi_1z_1\ne\bar\pi_2z_2$. For a uniformly random table $P'$ on $O$, what is the probability that the encodings of $z_1$ and $z_2$ give the shared vertex of $P'$ the same label?
\end{question}
